\chapter{Case Studies II}\label{ch:fm:case-studies-ii}

On 1 August 2012 a market maker, Knight Capital, lost more than \$460 million in forty-five minutes. Five days later it had sold \$400 million of
convertible preferred stock to a group of investors, convertible into about 73\% of its shares; eleven months later it merged with GETCO under a new
holding company. The controls failed in the morning; the decisions that set what the failure cost the firm's owners came in the next five days, and
many years before. This chapter reads four more failures from the public record: a software incident, a family office's prime brokers, an
exchange that was also a counterparty, and an exchange's decision in a squeeze.

\section{2012: a software incident and a rescue}

The Securities and Exchange Commission's order of October 2013 sets out the morning: a technician did not copy a new release of the firm's order router to
one of its eight servers, where an old function remained present and callable, and over 45 minutes it sent millions of orders and obtained more than 4 million executions in 154 stocks,
for more than 397 million shares. Book 11 read the order for its controls: no automated limit on the firm's aggregate exposure, and monitoring that
did not alert on it. The firm lost over \$460 million on the order's figures; its quarterly report put the realised pre-tax loss at about \$440 million,
against total equity of \$1\,497 million on 30 June.

\begin{dated}{2026-09}{The public record of the chapter's cases}\label{dat:fm:case-studies-ii:record}
\textbf{Knight Capital} (SEC order 34-70694; Forms 10-Q, 8-K of 6 August 2012 and KCG's 8-K12G3 of 1 July 2013): over \$460 million lost on 1
August 2012; \$400 million of convertible preferred stock sold on 6 August, convertible into about 73\% of the common stock, under the stock exchange's
financial viability exception to shareholder approval; merged with GETCO under KCG Holdings on 1 July 2013, each share exchanged for \$3.75 in cash
or a third of a KCG share. \textbf{Archegos} (Credit Suisse special committee report, 29 July 2021): Credit Suisse lost about \$5.5 billion; the Federal
Reserve Board fined UBS Group \$268.5 million for Credit Suisse's counterparty risk management on 24 July 2023, about \$387 million with the Prudential
Regulation Authority's penalty; the founder was convicted in July 2024 and sentenced to 18 years in prison. \textbf{LME nickel}: the Divisional Court
held that the exchange acted lawfully in cancelling the 8 March 2022 trades, and the Court of Appeal dismissed the appeal on 7 October 2024; the FCA
fined the exchange \pounds9\,245\,900 on 19 March 2025 for failing to maintain orderly trading.
\end{dated}

On 6 August the firm sold 400\,000 shares of convertible preferred stock for \$400 million to a group including Jefferies, Blackstone and GETCO,
convertible into about 266.7 million common shares, \$1.50 each, against 97.9 million shares outstanding. Issuing that much stock normally needs a
shareholder vote; the board's finance and audit committee relied instead on the exchange's exception for a delay that would seriously jeopardise the
company's financial viability, and the exchange did not object. Eleven months later the firm merged with GETCO under a new holding company, and each
share became \$3.75 in cash or a third of a share of the new company.

\begin{table}[ht]
\centering
\begin{tabular}{lrrr}
\toprule
 & 30 June 2012 & after the loss & after the rescue \\
\midrule
Book equity (\$ million) & 1\,497 & 1\,057 & 1\,457 \\
Shares (million) & 97.9 & 97.9 & 364.6 \\
Old shareholders' fraction & 100\% & 100\% & 26.9\% \\
Old shareholders' book equity (\$ million) & 1\,497 & 1\,057 & 391 \\
Book value per share (\$) & 15.29 & 10.80 & 4.00 \\
\bottomrule
\end{tabular}
\caption{The 2012 rescue on book values from the quarterly report and the purchase agreement, before tax: the loss took 29.4\% of the equity, the
rescue at \$1.50 a share took most of what remained from the old shareholders. Data: \texttt{fm\_cases2.knight}.}\label{tab:fm:case-studies-ii:rescue}
\end{table}

The table is the firm-level arithmetic of a rescue under time pressure. The loss alone cost the shareholders 29.4\% of the book equity; the rescue,
priced at \$1.50 a share when book value per share was still \$10.80, cost them a further \$666 million of book equity, transferred to the new
investors, who paid \$400 million for \$1\,066 million. The price was set by the days the firm had left: its clients and counterparties, the
quarterly report says, had reduced order flow and lost confidence, and a firm without capital cannot keep making markets. What it had not decided in
advance (the limit on aggregate exposure, the capital held against a morning like that one, a standing facility that could be drawn in a day) set the
terms.

\begin{omfigure}
\begin{tikzpicture}
\begin{axis}[width=9.5cm,height=5.2cm,ybar stacked,bar width=26pt,ymin=0,ymax=1700,ylabel={\small book equity (\$ million)},xtick={0,1,2},
  xticklabels={30 June 2012,after the loss,after the rescue},xmin=-0.6,xmax=2.6,tick label style={font=\footnotesize},grid=major,
  grid style={gray!20},table/col sep=comma,legend style={font=\scriptsize,at={(0.5,-0.2)},anchor=north,
  /tikz/every even column/.append style={column sep=8pt}},legend columns=2,legend cell align=left]
\addplot[fill=omDef!60,draw=omDef] table[x=pos,y=old]{figdata/desk/29-case-studies-ii/knight.csv};
\addplot[fill=omThm!45,draw=omThm] table[x=pos,y=new]{figdata/desk/29-case-studies-ii/knight.csv};
\coordinate (k0) at (axis cs:0,1497);
\coordinate (k1) at (axis cs:1,1057);
\coordinate (k2) at (axis cs:2,196);
\coordinate (k3) at (axis cs:2,924);
\legend{old shareholders,rescue investors}
\end{axis}
\node[font=\scriptsize,above] at (k0) {1\,497};
\node[font=\scriptsize,above] at (k1) {1\,057};
\node[font=\scriptsize] at (k2) {391};
\node[font=\scriptsize] at (k3) {1\,066};
\end{tikzpicture}
\omcaption{Who owned the 2012 market maker's book equity: the loss, then the rescue at \$1.50 a share. Data: \texttt{fm\_cases2.knight}.}\label{fig:fm:case-studies-ii:knight}
\end{omfigure}

\section{2021: a family office and its prime brokers}

The Credit Suisse board's special committee published its report on the bank's losses from Archegos, a family office, in July 2021. On the evening of
25 March 2021, after the value of its largest positions had fallen steeply since the 22nd, Archegos told its prime brokers that it had \$120 billion of gross
exposure, \$70 billion long and \$50 billion short, against \$9 to \$10 billion of equity, and asked them to agree a standstill while it wound its
positions down. They declined. Credit Suisse's own gross exposure was about \$27 billion on 23 March and \$17 billion on 26 March; it had called \$2.8
billion of margin on the 25th; it lost about \$5.5 billion.

\begin{definition}[Dynamic margining]\label{def:fm:case-studies-ii:dynamic}
\emph{Dynamic margining}\index{dynamic margining} sets a client's margin from the current risk of its portfolio (its volatility, concentration,
liquidity and directional bias), recomputed as positions and prices change, rather than as a fixed percentage of each trade's notional at inception.
\end{definition}

The report describes both regimes inside one bank. The prime brokerage book was dynamically margined; the swaps were margined statically, the initial
margin fixed in dollars at inception, so that as the positions rose in value the margin as a share of them eroded. The average margin on the swaps was
about 9.4\% on 23 March. A concentrated position that its own holder estimated would take two weeks to a month to liquidate carried the margin of a
diversified one.

\begin{definition}[Exit race]\label{def:fm:case-studies-ii:race}
An \emph{exit race}\index{exit race} is the sequential liquidation, by several creditors, of collateral or hedges in the same concentrated positions
after a common client defaults: each sells into prices moved by those who sold before it, so each creditor's loss depends on its place in the order as
well as on its margin.
\end{definition}

The report records the race. On 26 March Goldman organised block sales of some of the positions, and Credit Suisse took part in three without knowing
how many shares were offered or whose were sold first; it sold just over \$3 billion that day, keeping its algorithmic selling within 2 to 3\% of
average daily volume. Over the weekend several banks, including Deutsche Bank, Morgan Stanley and Goldman, were not interested in a managed
liquidation; Credit Suisse, UBS and Nomura agreed one and sold blocks in April.

\section{Tutorial: the rescue, the race, the run, the squeeze}\label{tut:fm:case-studies-ii:tut}

\begin{tutorial}
\textbf{Goal.} Extend the casebook to the four cases: the 2012 rescue's dilution; an \omterm{def:fm:case-studies-ii:race}{exit race} among seven prime brokers calibrated to the report;
the 2022 exchange's shortfall and run; and a nickel short's margin call against its cash. \textbf{End state:}
\cref{tab:fm:case-studies-ii:rescue} and \cref{fig:fm:case-studies-ii:knight,fig:fm:case-studies-ii:race,fig:fm:case-studies-ii:ftx,fig:fm:case-studies-ii:nickel}.
\begin{tutsteps}
\item \textbf{The rescue.} \texttt{firm.casebook.rescue} with the equity, the loss, \$400 million for 266.7 million shares, and 97.9 million old shares.
\item \textbf{The race.} Seven brokers (Credit Suisse and the six the report names) each hold one seventh of the position, each \$17 billion like Credit
Suisse on 26 March, about the \$120 billion gross in all. They sell in turn; the price falls linearly with the cumulative fraction sold (Book 10's
permanent impact).
\item \textbf{The calibration.} The impact coefficient is set so that a broker with Credit Suisse's exposure and its 9.4\% margin, selling sixth, loses
\$5.5 billion: \texttt{race\_impact} gives 0.531, a fall of 53\% by the time all seven have sold. The place is an assumption: the report puts Credit
Suisse among the brokers that sold into April.
\item \textbf{The run.} \texttt{ftx\_balances} and \texttt{ftx\_flows} read Book 3's transcription of the debtors' shortfall table.
\item \textbf{The squeeze.} An illustrative short of 10\,000 tonnes of nickel with \$300 million of cash, margined from the 7 March close.
\end{tutsteps}
\end{tutorial}

\omcode{code/firm/casebook/firm_casebook.py}{101}{116}{The exit race: brokers sell in turn under linear permanent impact, each losing the fall over its
segment net of its margin; and the impact coefficient that reproduces a stated loss.}\label{lst:fm:case-studies-ii:race}

\begin{omfigure}
\begin{tikzpicture}
\begin{axis}[width=10.5cm,height=5.6cm,xlabel={\small place in the exit race},ylabel={\small broker's loss (\$ billion)},xmin=0.7,xmax=7.3,
  ymin=-0.3,ymax=7.8,xtick={1,2,3,4,5,6,7},tick label style={font=\footnotesize},grid=major,grid style={gray!20},table/col sep=comma,
  legend cell align=left,legend style={font=\scriptsize,at={(0.03,0.97)},anchor=north west}]
\addplot[omThm,thick,mark=triangle*] table[x=place,y=m075]{figdata/desk/29-case-studies-ii/race.csv};
\addplot[omDef,very thick,mark=*] table[x=place,y=m094]{figdata/desk/29-case-studies-ii/race.csv};
\addplot[gray,thick,dashed,mark=square*] table[x=place,y=m200]{figdata/desk/29-case-studies-ii/race.csv};
\draw[black,thin] (axis cs:6,5.5) circle[radius=3pt];
\node[font=\scriptsize,anchor=south east] at (axis cs:5.9,5.7) {calibration};
\legend{margin 7.5\%,margin 9.4\%,margin 20\%}
\end{axis}
\end{tikzpicture}
\omcaption{A stylised \omterm{def:fm:case-studies-ii:race}{exit race} among seven prime brokers, each with \$17 billion of exposure to the same positions, calibrated so that the sixth,
with a 9.4\% margin, loses \$5.5 billion. The first seller loses nothing; the last loses \$6.79 billion at 9.4\% and \$4.99 billion at 20\%. Data:
\texttt{fm\_cases2.race}.}\label{fig:fm:case-studies-ii:race}
\end{omfigure}

In the model the first broker to sell loses nothing at any of the three margins: it sells at an average fall of 3.8\%, inside its margin. The last
loses \$6.79 billion at 9.4\%, and \$4.99 billion even at 20\%. The margin that would have left each broker whole is the average fall over its
segment: 3.8\% for the first, 26.6\% for the middle, 49.3\% for the last. No static percentage covers the race, because the loss depends on the
order, which no broker controls in advance.

\begin{proposition}[The race redistributes a fixed fall]\label{prop:fm:case-studies-ii:zero}
With linear permanent impact $1-\kappa f$ and brokers holding shares $s_1,\dots,s_K$ of the position, summing to one, the total fall in value
realised by all the brokers is $\kappa/2$ of the position, whatever the order; the order only decides who bears it. The broker in place $k$ bears
$s_k\kappa\bigl(F_{k-1}+s_k/2\bigr)$, where $F_{k-1}$ is the share sold before it.
\end{proposition}

\begin{proof}
Broker $k$ sells between the cumulative fractions $F_{k-1}$ and $F_k=F_{k-1}+s_k$ at the average price $1-\kappa(F_{k-1}+s_k/2)$, so its fall is
$s_k\kappa(F_{k-1}+s_k/2)=\kappa\int_{F_{k-1}}^{F_k}f\,df$. Summing over $k$ gives $\kappa\int_0^1f\,df=\kappa/2$.
\end{proof}

The proposition is why a standstill is rational for the creditors together and irrational for each one alone: a managed liquidation leaves the total
fall unchanged (or smaller, if slower selling lowers the impact), while selling first moves one's own share of it onto the others. \omterm{def:fm:case-studies-ii:dynamic}{Dynamic margining}
is the defence each broker controls: a margin that grows with the position's size against the market's volume and with the broker's ignorance of
what the client holds elsewhere, set before the race starts.

\section{2022: an exchange that was also a counterparty}

Book 3 read FTX's collapse from the debtors' filings. For a trading firm the lesson is about exposure. A firm that holds cash and tokens on an exchange
is a creditor of the exchange; if the exchange lends customers' assets to an affiliate, the firm is also, without knowing it, a creditor of that
affiliate. The debtors' first interim report found that the affiliate, Alameda, had a borrowing limit set to \$65 billion and settings that let it
withdraw below it and run a negative balance; the shortfall analysis found that at the petition, for the most liquid tokens, the exchange held assets
worth 6.6\% of what it owed customers, and 22.6\% for all tokens with receivables.

\begin{omfigure}
\begin{tikzpicture}
\begin{axis}[width=10.5cm,height=5.4cm,xlabel={\small November 2022},ylabel={\small cumulative net flow (\$ billion)},xmin=1,xmax=11,ymin=-7.8,ymax=4,
  xtick={1,2,3,4,5,6,7,8,9,10,11},tick label style={font=\footnotesize},grid=major,grid style={gray!20},table/col sep=comma,legend cell align=left,
  legend style={font=\scriptsize,at={(0.03,0.97)},anchor=north west}]
\addplot[omDef,very thick,mark=*,mark size=1.5pt] table[x=day,y=customers]{figdata/desk/29-case-studies-ii/ftx.csv};
\addplot[omThm,thick,dashed,mark=square*,mark size=1.5pt] table[x=day,y=related]{figdata/desk/29-case-studies-ii/ftx.csv};
\draw[gray,dotted] (axis cs:8,-7.8) -- (axis cs:8,4) node[pos=0.3,right,font=\scriptsize] {withdrawals paused};
\draw[black,thin] (axis cs:1,0) -- (axis cs:11,0);
\legend{customers,Alameda and related parties}
\end{axis}
\end{tikzpicture}
\omcaption{The run on the 2022 exchange: customers' cumulative net withdrawals reached \$7.0 billion by the petition, \$3.3 billion of it on 7 November
alone, while the affiliate and related parties deposited a net \$3.2 billion. Data: \texttt{fm\_cases2.ftx}, from Book 3's transcription of the
debtors' table.}\label{fig:fm:case-studies-ii:ftx}
\end{omfigure}

The run was fast: customers withdrew a net \$3.3 billion on 7 November, and on 8 November the exchange paused withdrawals. A firm that learned of the
shortfall on the 8th had no exit left; the decisions that mattered were the limits set earlier on how much the firm held on any one venue, how
quickly excess balances were swept back, and what evidence of segregation it required (chapter 27's \omterm{def:fm:crisis-management:trapped}{trapped assets}).

\section{2022: the nickel squeeze and the exchange's decision}

On 7 March 2022 the London Metal Exchange's three-month nickel contract closed at \$48\,078 a tonne, up from an opening just under \$30\,000. In
the first hours of 8 March it rose to \$101\,365 at 06:08 and stayed above \$80\,000 from 07:00. The Court of Appeal's judgment records what the
exchange faced: if the morning's trades stood, members would have had to post some \$19.75 billion of margin by 09:00, and at least five were
expected to default. It suspended trading at 08:15 and cancelled every trade made since midnight. Large short positions had been built by several
participants, among them Tsingshan on the over-the-counter market. The FCA's notice adds that margin calls had set records of \$3.5 billion on 4
March and \$5.1 billion on the morning of 7 March, and that the price bands were switched off from 04:49 to 08:15 on a decision by junior staff.

\begin{omfigure}
\begin{tikzpicture}
\begin{axis}[width=10.5cm,height=5.4cm,xlabel={\small nickel price (\$ thousand a tonne)},ylabel={\small variation margin call (\$ million)},
  xmin=46,xmax=104,ymin=0,ymax=600,tick label style={font=\footnotesize},grid=major,grid style={gray!20},table/col sep=comma]
\addplot[omDef,very thick] table[x=price,y=call]{figdata/desk/29-case-studies-ii/nickel.csv};
\draw[omThm,thick,dashed] (axis cs:46,300) -- (axis cs:104,300) node[pos=0.02,above right,font=\scriptsize] {cash available};
\draw[gray,dotted] (axis cs:78.078,0) -- (axis cs:78.078,600) node[pos=0.93,left,font=\scriptsize] {cash runs out};
\draw[gray,dotted] (axis cs:101.365,0) -- (axis cs:101.365,600) node[pos=0.93,left,font=\scriptsize] {peak};
\end{axis}
\end{tikzpicture}
\omcaption{An illustrative short of 10\,000 tonnes of nickel margined from the 7 March close: its call reaches \$319 million at \$80\,000 and \$533
million at the peak, and its \$300 million of cash runs out at \$78\,078. Data: \texttt{fm\_cases2.nickel}.}\label{fig:fm:case-studies-ii:nickel}
\end{omfigure}

For the short, the arithmetic is chapter 14's: a call that grows linearly with the price, against cash that does not. For the exchange it was a
choice between a margin call that would have defaulted several members and a cancellation that would take trades from those who made them. The
Divisional Court held the cancellation lawful, and dismissed Jane Street's similar claim; the Court of Appeal dismissed Elliott's appeal, holding that
the only alternative, collecting margin at the previous close, would have put the clearing house in breach of its own obligations. Two firm-level
lessons follow for a trading firm: an exchange can cancel trades in extreme conditions, so a profit made in a disorderly market is not yet a profit;
and the firm's margin plan must survive calls several times larger than any seen before, since the record ones came on consecutive trading days.

\section{What the four share}

\begin{method}[From a case to a control]\label{met:fm:case-studies-ii:control}
\begin{enumerate}
\item State the decision the record shows was taken, or not taken, before the event: a limit, a margin method, a venue exposure, a funding plan.
\item Size the event on the public figures: the loss against equity, the race's losses by place, the shortfall, the call against cash.
\item Write the control that would have changed the size: an aggregate exposure limit and a capital plan for it; dynamic margins with concentration
add-ons; limits and sweeps on venue balances; a margin plan for record calls.
\item Test the control on the firm's own book with the casebook's model, and record the result in the risk register (chapter 12).
\end{enumerate}
\end{method}

The four failures differ in cause and share a structure. In each, a counterparty's terms were fixed before the event on assumptions the event broke: a
rescue priced in days, a margin set at inception, a balance left on a venue, a margin call sized for ordinary moves. And in each the firms that did
best had decided their response in advance: the brokers that sold first, the exchange with the power to cancel, the firms whose exposures were already
small. Chapter 27's drill is where a firm finds out which of these decisions it has taken.

\section{Build: the casebook, part 2}\label{bld:fm:case-studies-ii:casebook}

\begin{build}
\bfield{Purpose} The four cases' public figures as tables tied to their ledger rows; a rescue-dilution calculator; the \omterm{def:fm:case-studies-ii:race}{exit race} with a calibration
to a stated loss; the exchange shortfall and run from Book 3's data; a short's margin call against its cash.
\bfield{Interface} \texttt{firm.casebook}: \texttt{KNIGHT}, \texttt{ARCHEGOS}, \texttt{NICKEL}, \texttt{rescue}, \texttt{exit\_race},
\texttt{race\_impact}, \texttt{ftx\_balances}, \texttt{ftx\_flows}, \texttt{margin\_call}, \texttt{break\_price}.
\bfield{Rules} Figures come from the ledger or from committed data with a licence row; illustrative inputs are named as such; the race's total fall is
independent of the order (\cref{prop:fm:case-studies-ii:zero}).
\bfield{Acceptance tests} \texttt{code/firm/casebook/tests/}: the race's zero-sum property and its calibration, the rescue on hand numbers, the margin
call and break price, and the shortfall table's coverage.
\bfield{Stretch} Square-root impact in the race with each broker's own selling speed; a managed liquidation as a joint schedule; the nickel short
with initial margin raised by the clearing house.
\end{build}

\begin{omsources}
\item US Securities and Exchange Commission, In the Matter of Knight Capital Americas LLC, Release 34-70694, 16 October 2013.
\item Knight Capital Group, Form 10-Q for the quarter ended 30 June 2012; Form 8-K, 6 August 2012; KCG Holdings, Form 8-K12G3, 1 July 2013.
\item Credit Suisse Group Special Committee of the Board of Directors, \emph{Report on Archegos Capital Management}, 29 July 2021.
\item Board of Governors of the Federal Reserve System, press release, 24 July 2023; US Attorney's Office, SDNY, press release, 19 December 2024.
\item FTX Debtors, \emph{Preliminary Analysis of Shortfalls}, 2 March 2023; J.~J. Ray III, First Interim Report, 9 April 2023.
\item R (Elliott Associates) v London Metal Exchange [2024] EWCA Civ 1168; Financial Conduct Authority, Final Notice to the London Metal Exchange,
19 March 2025.
\end{omsources}

\section{Exercises}

\begin{exercise}[$\star$]\label{exo:fm:case-studies-ii:1}
From \cref{tab:fm:case-studies-ii:rescue}, what fraction of the 2012 firm did its old shareholders own after the rescue, and what price per share
did the rescue investors pay?
\end{exercise}

\begin{exercise}[$\star$]\label{exo:fm:case-studies-ii:2}
Define \omterm{def:fm:case-studies-ii:dynamic}{dynamic margining} and say why static margin on a rising concentrated position erodes.
\end{exercise}

\begin{exercise}[$\star$]\label{exo:fm:case-studies-ii:3}
A short of 10\,000 tonnes is margined from \$48\,078. What is its call at \$80\,000, and at what price does \$300 million of cash run out?
\end{exercise}

\begin{exercise}[$\star\star$]\label{exo:fm:case-studies-ii:4}
Prove \cref{prop:fm:case-studies-ii:zero}. What does it imply for a standstill agreement among creditors?
\end{exercise}

\begin{exercise}[$\star\star$]\label{exo:fm:case-studies-ii:5}
In the calibrated race, what margin leaves the first, the sixth and the last broker whole, and what margin covers the position if one broker held
all of it?
\end{exercise}

\begin{exercise}[$\star\star$]\label{exo:fm:case-studies-ii:6}
Which limits would have protected a trading firm with balances on the 2022 exchange, and which records would have told it the balances were at risk?
\end{exercise}

\begin{exercise}[$\star\star\star$]\label{exo:fm:case-studies-ii:7}
\emph{Coding.} Rerun the race with a 20\% margin. Which places lose nothing, and what does the last lose?
\end{exercise}

\begin{exercise}[$\star\star\star$]\label{exo:fm:case-studies-ii:8}
\emph{Find the flaw.} ``Our margin on the client is 20\%, twice the industry's, so we cannot lose on it.''
\end{exercise}

\section{Problem: The Last to Sell}

\begin{problem}[{Weekend problem --- the last to sell}]\label{pb:fm:case-studies-ii:1}
The head of a prime brokerage business must explain to the board what the four cases mean for its margin policy and its own exposures.

\medskip\noindent\textbf{Part I --- The concepts.}
\begin{enumerate}
\item Define \omterm{def:fm:case-studies-ii:dynamic}{dynamic margining}.
\item Define an \omterm{def:fm:case-studies-ii:race}{exit race}.
\item State \cref{prop:fm:case-studies-ii:zero}.
\item Why is a standstill rational for creditors together and not for each alone?
\end{enumerate}

\medskip\noindent\textbf{Part II --- The record.}
\begin{enumerate}[resume]
\item Summarise the public record of the 2012 incident and rescue.
\item What did the family office tell its prime brokers on 25 March 2021, and what did they do?
\item Summarise the 2022 exchange's shortfall and run.
\item What did the nickel market's exchange decide, and what did the courts and the FCA find?
\end{enumerate}

\medskip\noindent\textbf{Part III --- The reconstructions.}
\begin{enumerate}[resume]
\item Give the rescue's dilution table.
\item Describe the \omterm{def:fm:case-studies-ii:race}{exit race} model and its calibration.
\item Give the losses by place at a 9.4\% margin.
\item Give the losses by place at 20\%, and the margin each place needs.
\item What does the model leave out?
\item Give the illustrative nickel short's calls and break price.
\end{enumerate}

\medskip\noindent\textbf{Part IV --- The policy.}
\begin{enumerate}[resume]
\item What margin method would you adopt for concentrated clients?
\item How would you learn what a client holds at other brokers?
\item What limits would you set on the firm's own balances at venues?
\item What would you change in the firm's capital plan after the 2012 case?
\item State the \emph{named result}: each prime broker's loss as a function of its place in the \omterm{def:fm:case-studies-ii:race}{exit race} and of its margin, calibrated to the public
figures, and the dilution of the 2012 rescue.
\item In two sentences, write the recommendation.
\end{enumerate}
\end{problem}

\section{Interview questions}

\begin{interviewq}[$\star$ \iqroles{risk}]\label{iq:fm:case-studies-ii:1}
What is the difference between static and \omterm{def:fm:case-studies-ii:dynamic}{dynamic margining}?
\end{interviewq}

\begin{interviewq}[$\star$ \iqroles{developer, risk}]\label{iq:fm:case-studies-ii:2}
A new release starts sending orders nobody intended. Which controls should have stopped it, and which decision determines what it costs the owners?
\end{interviewq}

\begin{interviewq}[$\star\star$ \iqroles{risk, trader}]\label{iq:fm:case-studies-ii:3}
A client defaults on swaps you and five other banks hold on the same stocks. The others are proposing a standstill. What do you do?
\end{interviewq}

\begin{interviewq}[$\star\star$ \iqroles{risk}]\label{iq:fm:case-studies-ii:4}
How much of the firm's cash would you hold on a single crypto exchange, and why?
\end{interviewq}

\begin{interviewq}[$\star\star$ \iqroles{trader}]\label{iq:fm:case-studies-ii:5}
You made a large profit on a trade in a market that the exchange then declared disorderly. What should you expect?
\end{interviewq}

\begin{interviewq}[$\star\star\star$ \iqroles{researcher, risk}]\label{iq:fm:case-studies-ii:6}
Model the losses of several creditors liquidating the same collateral in turn. What does the order change, and what does it not?
\end{interviewq}
